\section{タスクの設定とデータセットの要件} \label{sec3}
本章では，家庭環境の理解を必要とする家庭内タスクについて説明し，その評価に必要なデータセットの概要を述べる．家庭環境知識とは，家庭環境内の物体に関する詳細な情報を指し，物体の現在の状態，識別番号，および物体間の関係を含む．これらの家庭環境知識に依存する家庭内タスクを評価するために必要なデータセットの要件を以下に示す．

\begin{itemize}
    \item 物体の識別番号の特定が必要なタスクとすること．これは，同じ種類の物体が複数存在する場合に，どの物体を対象とすべきかを明確にするためである．
    \item 物体の初期状態を常識的に推論すると失敗するタスクを含めること．例えば，``Turn on the light''という指示を受けた場合，常識的にはライトの状態が``OFF''であると考えられる．しかし，実際には``ON''である可能性もある．本研究では，常識知識のみに依存せず，家庭環境知識を活用することを目的とするため，このようなタスクを含める．
    \item タスクの達成条件（ゴール）が一意に定まるタスクとすること．家庭内タスクの例として``Clean up''が挙げられる．しかし，このようなタスクはゴールの定義が曖昧であり，解釈が複数存在するため，自動評価が困難である．通常は人手による評価が必要となるが，これには高い評価コストが伴うため，ゴールが明確に定義されたタスクとする．
    \item シミュレータ内部の環境グラフを取得可能とし，視覚情報による環境認識は行わないこと．その変換と，得られた記述を用いる計画過程を評価対象とする．
\end{itemize}

本研究では，上記の要件を統制して評価するため，同一種類の物体の識別と，初期状態・配置関係に応じた行動の変更を必要とするデータセットを構成する．評価はシミュレータ内部の環境グラフを利用する条件で行う．また，ゴール条件をデータとして保持し，指定された物体IDの状態または関係に基づいて評価する．

そこで，本研究では，評価に必要な要件を満たす2種類のタスクを設定した．一つ目は，複数の物体を特定の状態に変化させる「状態変化タスク」であり，主に物体の状態を認識することが課題となるタスクである．もう一つは，複数の物体を特定の場所（物体の上または中）に配置する「配置タスク」であり，主に物体間の関係を認識することが課題となるタスクである．

\subsection{指示文の範囲と粒度}\label{sec3.instructions}
本論文では，自然言語の指示文を「タスク名」，指示文とシーン・初期条件・ゴール条件・参照行動列を組にしたデータ1件を「タスクインスタンス」と呼ぶ．評価対象は，状態変化では``Turn on/off all [対象クラス]''，配置では``Put all [対象クラス] in/on the [配置先クラス]''という定型の指示である．各インスタンスは一つの対象クラスを扱い，配置タスクでは一つの配置先を指定する．

``all''は，当該シーン（家屋）内でタスクが対象とするクラスの物体全体を指す設定である．初期部屋にある物体だけに限定する指示ではない．対象ID，個数，現在状態，所属部屋は指示文に列挙せず，環境知識から取得する．例えば，``Turn on all lightswitches in the kitchen''は対象範囲を限定し，``Turn on all 4 lightswitches''は個数を追加する別条件であり，今回の指示文集合には含まれない．ただし，公開実装のプロンプトには，この``all''の規約を別途説明する追加文はない．実装はタスク中の名詞に対応する物体を家屋全体のグラフから抽出する．

\subsection{入力，初期化および評価情報}\label{sec3.inputs}
データの各フィールドの用途を\ref{tab:input_roles}に示す．エージェントの初期位置は\texttt{initial\_room}で指定し，実行中の所属部屋・近接関係はグラフから取得する．移動は対象名とIDを指定した高レベル操作WALKによる．したがって，初期部屋は初期化と実行状態に関係するが，物体検索の範囲をその部屋に制限するものではない．経路上の前進・回転等を本手法が逐次計画する設定ではない．

\begin{table*}[tb]
\caption{データの入力先と用途（公開実装）}\label{tab:input_roles}
\centering\small
\begin{tabular}{p{.23\textwidth}p{.69\textwidth}}
\hline
フィールド／情報 & 用途 \\
\hline
\texttt{task} & 計画・完了判定の指示，関連名詞の抽出，事例検索のクエリ．\\
\texttt{scene}, \texttt{initial\_room} & シーンを初期化し，指定部屋にキャラクターを配置する．これらのフィールドをそのままプロンプトに渡す処理はない．\\
\texttt{initial\_states} & 指定IDのノードの状態配列を上書きする．物体位置を変更するフィールドではない．\\
実行中の環境グラフ & \ref{sec4.1}の抽出・自然言語化を経て，計画・完了判定・前提条件判定に渡す．\\
\texttt{goal\_states} & 通常終了時の採点に利用し，計画用プロンプトには渡さない．\\
評価対象の\texttt{action\_scripts} & 参照行動列の長さを試行上限の設定に利用する．行動列そのものは計画器に渡さない．\\
事例用データの\texttt{task}, \texttt{action\_scripts} & 選択した1件の指示・行動列をプロンプトに埋め込む．事例の初期状態・シーンは埋め込まない．\\
\hline
\end{tabular}
\end{table*}

\subsection{推論の範囲とゴール以外の制約}\label{sec3.scope}
評価する推論は，同名物体の識別，現在状態に応じた操作対象の選択，物体の支持・包含関係の把握，配置前の把持や容器開放といった局所的な操作条件の判断である．複雑さを表す属性として，ゴール対象数，状態変更が必要な対象数，配置先の開閉状態，シーン，初期部屋，参照行動列長を記録する．

タスクの構築方針として，目標遂行のために障害物を移動させる操作や，達成済みのゴールを一時的に崩す後戻りを行わずに達成できるタスクを対象とする．必要な物体の状態や関係を把握し，把持・移動・容器開放・配置などを適切な順序で行う推論を評価する．

ゴール条件は最終状態について定義する．配置タスクでは指定ID間のONまたはINSIDE関係を評価し，指示にない片付け，無関係な物体の原状回復，冷蔵庫の最終的な閉鎖や開放時間は独立の採点条件に含めない．途中状態の制約としてはシミュレータ自身の実行条件を適用し，物体を一つ格納するごとに冷蔵庫を閉め直すことは要求しない．参照行動列は実行方法の一例である．実行エラーと試行上限を含む成功・失敗の判定順序は\ref{sec5.2}と\ref{sec5.3}に示す．
